% ============================================================
\part{Docker}
% ============================================================

\chapter{Wir wollen ein Programm auf jedem Computer gleich ausführen}

\kapitelziel{
Du verstehst das Container-Konzept und startest einen
ROS2-Container, in dem \texttt{stufe2\_ros2/} läuft --
identisch auf Laptop, Raspberry Pi und Schulrechner.
}
\kapitelstruktur

\section{Situation}

\texttt{sensor\_bridge.py} läuft auf deinem Laptop.
Auf dem Raspberry Pi des Roboters fehlt \texttt{pyserial}.
Auf dem Schulrechner ist Ubuntu 20.04, aber du entwickelst mit 24.04.
Jedes Mal andere Umgebung, jedes Mal andere Fehler.

\begin{aufgabeBox}
Starte \texttt{sensor\_bridge.py} in einem Docker-Container,
der auf jedem Linux-System identisch funktioniert.
\end{aufgabeBox}

\begin{problemBox}
Software hat Abhängigkeiten. Diese haben Abhängigkeiten.
Auf dem eigenen Rechner funktioniert alles -- woanders nicht.
Das nennt sich \textit{Dependency Hell}.
Docker ist die Lösung.
\end{problemBox}

\begin{wissenBox}
Wir brauchen: den Unterschied zwischen Image und Container,
grundlegende Docker-Befehle und das Prinzip der Isolation.
\end{wissenBox}

\begin{haubeBox}
\textbf{Warum ist Docker kein Betrug an der Portierbarkeit?}

Früher gab es VMs: Ein vollständiger zweiter Betriebssystem-Kernel
lief innerhalb des Hosts. Teuer, langsam, gigantisch.

Docker nutzt Linux-Kernel-Funktionen -- \textit{namespaces}
für Isolation und \textit{cgroups} für Ressourcenbegrenzung.
Der Container läuft im selben Kernel wie der Host.
Das Ergebnis: Ein ROS2-Container startet in Sekunden statt Minuten.
Und er ist auf jedem Linux identisch -- weil Linux gleich ist.
\end{haubeBox}

\begin{analogie}
Ein Docker-Image ist wie die ISO-Datei einer Linux-Installation.
Ein Container ist das laufende System aus diesem Image --
wie eine gestartete VM, aber ohne eigenen Kernel.
Aus einem Image können viele Container entstehen.
\end{analogie}

\section{Der pragmatische Weg}

\begin{lstlisting}[style=bash]
# Docker einrichten (einmalig)
sudo apt install docker.io
sudo usermod -aG docker $USER
newgrp docker

# ROS2-Image laden
docker pull ros:jazzy

# Erster Test: ROS2 im Container starten
docker run -it ros:jazzy bash
# Jetzt sind wir INSIDE des Containers:
source /opt/ros/jazzy/setup.bash
ros2 --version
exit
\end{lstlisting}

\begin{loesungBox}
\begin{lstlisting}[style=bash]
# Container mit Projektordner starten
docker run -it \
  -v ~/robotik-uno-q:/robotik-uno-q \
  ros:jazzy bash

# Im Container:
cd /robotik-uno-q/stufe2_ros2
pip3 install pyserial pyyaml
python3 src/sensor_bridge.py
\end{lstlisting}
Pyserial ist jetzt im Container -- nicht am Host.
Das Projektverzeichnis ist per Volume eingebunden:
Änderungen am Host sind sofort im Container sichtbar.
\end{loesungBox}

\begin{projektBox}
\textbf{Bisher:} \texttt{sensor\_bridge.py} läuft nur lokal.\\
\textbf{Heute:} ROS2-Container gestartet, Projekt eingebunden.
Läuft auf jedem Linux.\\
\textbf{Nächstes Kapitel:} Im Container debuggen.
\end{projektBox}

\begin{merksatzBox}
Image = Bauplan (unveränderlich).
Container = laufende Instanz.
\texttt{-v host:container} bindet Ordner ein.
Container sind kurzlebig -- Daten überleben nur via Volumes.
\end{merksatzBox}

\begin{robotikBox}
Das offizielle \texttt{ros:jazzy}-Image enthält eine vollständige
ROS2-Installation. In der Robotik-Industrie ist Docker Standard:
Jeder Entwickler klont das Repo und startet einen Container --
keine Installationsanleitung, keine Abhängigkeitsprobleme.
\end{robotikBox}

% ── Kapitel 17 ────────────────────────────────────────────────
\chapter{Wir wollen einen laufenden Container betreten}

\kapitelziel{
Du kannst in laufende Container einsteigen und
\texttt{sensor\_bridge.py} von innen debuggen.
}
\kapitelstruktur

\section{Situation}

\texttt{sensor\_bridge.py} läuft im Container, aber liefert
falsche Werte. Du musst von innen nachschauen:
Welche Umgebungsvariablen sind gesetzt? Welche Python-Pakete
sind installiert? Sieht der Container den Serial-Port?

\begin{aufgabeBox}
Starte den Container im Hintergrund und betrete ihn
mit einer zweiten Shell, um das Problem zu untersuchen.
\end{aufgabeBox}

\begin{lstlisting}[style=bash]
# Container im Hintergrund starten
docker run -d --name robotik-dev \
  -v ~/robotik-uno-q:/robotik-uno-q \
  ros:jazzy sleep infinity

# In laufenden Container einsteigen
docker exec -it robotik-dev bash

# Innerhalb des Containers:
pip3 list | grep serial   # Pyserial installiert?
ls /dev/ttyUSB*           # Serial-Port sichtbar?
env | grep ROS            # ROS2-Umgebung geladen?
python3 -c "import serial; print(serial.__version__)"
exit

# Logs des Containers anzeigen
docker logs robotik-dev
docker logs -f robotik-dev   # Live

# Aufraeumen
docker stop robotik-dev
docker rm robotik-dev
\end{lstlisting}

\begin{projektBox}
\textbf{Heute:} Laufenden Container betreten, Konfiguration
von innen geprüft.\\
\textbf{Wichtige Erkenntnis:} \texttt{/dev/ttyUSB0} ist im Container
nicht sichtbar -- das beheben wir in Kapitel~23.
\end{projektBox}

\begin{merksatzBox}
\texttt{docker exec -it NAME bash} öffnet eine Shell
in einem laufenden Container -- wie SSH, aber für Container.
Es ist das wichtigste Debugging-Werkzeug nach \texttt{docker run}.
\end{merksatzBox}

% ── Kapitel 18 ────────────────────────────────────────────────
\chapter{Wir wollen Daten behalten, obwohl der Container gelöscht wird}

\kapitelziel{
Du richtest Volumes ein, damit Sensor-Logs aus dem Container
dauerhaft in \texttt{stufe2\_ros2/logs/} gespeichert werden.
}
\kapitelstruktur

\section{Situation}

\texttt{sensor\_bridge.py} schreibt Logs.
Der Container wird nach dem Testen gelöscht.
Die Logs sind weg.

\begin{aufgabeBox}
Binde \texttt{stufe2\_ros2/logs/} als Volume ein, sodass
Sensor-Logs den Container-Neustart überleben.
\end{aufgabeBox}

\begin{lstlisting}[style=bash]
# Bind Mount: Host-Ordner im Container verfuegbar machen
docker run -it \
  -v ~/robotik-uno-q:/robotik-uno-q \
  -v ~/robotik-uno-q/stufe2_ros2/logs:/root/.ros/log \
  ros:jazzy bash

# Im Container: Logs schreiben
cd /robotik-uno-q/stufe2_ros2
python3 src/sensor_bridge.py 2>&1 | tee /root/.ros/log/session.log

exit

# Auf dem Host: Log ist noch da!
cat ~/robotik-uno-q/stufe2_ros2/logs/session.log
\end{lstlisting}

\begin{projektBox}
\textbf{Heute:} \texttt{stufe2\_ros2/logs/} als Bind Mount --
Logs überleben Container-Neustarts.
\end{projektBox}

\begin{robotikBox}
ROS2-Rosbag-Aufnahmen (\texttt{.db3}) werden als Bind Mount
gespeichert. So gehen Messdaten nie verloren, auch wenn der
Container für ein Update neu gebaut wird.
Das Bag bleibt, der Container kann weg.
\end{robotikBox}

% ── Kapitel 19 ────────────────────────────────────────────────
\chapter{Wir wollen unser eigenes Docker-Image bauen}

\kapitelziel{
Du erstellst ein Dockerfile für \texttt{stufe2\_ros2/}
mit allen Python-Abhängigkeiten von robotik-uno-q.
}
\kapitelstruktur

\section{Situation}

Jedes Mal \texttt{pip3 install pyserial pyyaml} im Container
tippen ist nervig. Das Dockerfile automatisiert das -- und
macht den Container reproduzierbar.

\begin{aufgabeBox}
Erstelle \texttt{stufe2\_ros2/Dockerfile} mit allen
Abhängigkeiten für \texttt{sensor\_bridge.py}.
\end{aufgabeBox}

\begin{lstlisting}[style=dockerfile]
# stufe2_ros2/Dockerfile
FROM ros:jazzy

LABEL maintainer="Wolfgang Becker"
LABEL description="robotik-uno-q Stufe 2 -- Sensor-Bridge"

# Python-Abhaengigkeiten installieren
RUN pip3 install pyserial pyyaml

# Arbeitsverzeichnis
WORKDIR /robotik-uno-q/stufe2_ros2

# Startbefehl (ueberschreibbar per docker run)
CMD ["python3", "src/sensor_bridge.py"]
\end{lstlisting}

\begin{lstlisting}[style=bash]
cd ~/robotik-uno-q/stufe2_ros2

# Image bauen
docker build -t robotik-uno-q:v0.1 .

# Testen
docker run -it \
  -v ~/robotik-uno-q:/robotik-uno-q \
  robotik-uno-q:v0.1 bash
\end{lstlisting}

\begin{projektBox}
\textbf{Heute:} \texttt{stufe2\_ros2/Dockerfile} erstellt.
Jeder, der das Repo klont, kann mit einem Befehl die
identische Umgebung aufbauen.
\end{projektBox}

\begin{merksatzBox}
Jede \texttt{RUN}-Zeile erzeugt einen Layer.
Wenige kombinierte Befehle = kleineres Image.
\texttt{FROM ros:jazzy} gibt es gratis -- keine manuelle
ROS2-Installation nötig.
\end{merksatzBox}

% ── Kapitel 20 ────────────────────────────────────────────────
\chapter{Wir wollen mehrere Container gemeinsam starten}

\kapitelziel{
Du definierst den vollständigen stufe2\_ros2-Stack
in \texttt{docker-compose.yml} und startest alles
mit einem einzigen Befehl.
}
\kapitelstruktur

\section{Situation}

Bald hat \texttt{stufe2\_ros2/} mehrere Komponenten:
die Serial-Bridge zum Arduino, die Navigation und
ein einfaches Dashboard. Drei separate
\texttt{docker run}-Befehle -- das ist fehleranfällig.

\begin{aufgabeBox}
Definiere den Stack in \texttt{docker-compose.yml}
und starte alle Dienste mit \texttt{docker-compose up}.
\end{aufgabeBox}

\begin{lstlisting}[style=yaml]
# robotik-uno-q/docker-compose.yml
version: '3.8'

services:
  sensor_bridge:
    build: ./stufe2_ros2
    volumes:
      - ./stufe2_ros2:/robotik-uno-q/stufe2_ros2
      - ./stufe2_ros2/logs:/root/.ros/log
    devices:
      - "/dev/arduino:/dev/arduino"   # Arduino-Port (nach Kapitel 23)
    restart: unless-stopped
    networks:
      - ros_net

  # Platzhalter fuer Band 2:
  # navigation_node:
  #   ...

networks:
  ros_net:
    driver: bridge
\end{lstlisting}

\begin{lstlisting}[style=bash]
cd ~/robotik-uno-q

# Stack starten
docker-compose up -d

# Logs verfolgen
docker-compose logs -f sensor_bridge

# Stoppen
docker-compose down
\end{lstlisting}

\begin{projektBox}
\textbf{Heute:} \texttt{docker-compose.yml} im Projektstamm --
der gesamte stufe2-Stack mit einem Befehl.\\
\textbf{Zusammen mit systemd aus Kapitel~10:} Nach dem Einschalten
startet systemd Docker, Docker startet den Stack.
Komplett automatisch.
\end{projektBox}

\begin{merksatzBox}
\texttt{docker-compose up -d} startet alles.
\texttt{docker-compose down} stoppt alles.
\texttt{devices:} gibt dem Container Zugriff auf
Hardware-Geräte -- ohne das sieht der Container
den Arduino nicht.
\end{merksatzBox}

\begin{robotikBox}
In professionellen Robotikprojekten definiert die
\texttt{docker-compose.yml} den kompletten Produktivstack:
Sensor-Nodes, Navigation, Logging, Monitoring.
Die Datei im Repo ist gleichzeitig Dokumentation
und Deploymentanleitung.
\end{robotikBox}
